\section{Auswertung des Jahresabschlusses}
\label{sec:abschluss-rdy}

Berichtender Abschnitt, jede Kennzahl mit Rechenweg. Alle Betr\"age in
Mio.\,INR. Das Gesch\"aftsjahr FY2026 endete am 31.~M\"arz 2026.

\subsection{Umsatz}

Der Umsatz stieg \Pct{\RdyUmsatzWachstum} auf \MioINR{\RdyUmsatzXXVI} nach
\MioINR{\RdyUmsatzXXV}.\quelleZF{2026}{112}\quelleMerken{s25zfxcacgxbbc}
Das Unternehmen weist selbst aus, dass dieser Zuwachs ein W\"ahrungseffekt
ist: W\"ahrungsbereinigt ging der Umsatz um \PctBetrag{\RdyOrganisch} zur\"uck, und
als Ursache nennt der Bericht Preissenkungen bestehender Produkte
einschlie{\ss}lich Lenalidomid in den Vereinigten
Staaten.\quelleZF{2026}{54}\quelleMerken{s25zfxcacgxfe}
\[
  \text{Umsatzwachstum} = \frac{\RdyUmsatzXXVI}{\RdyUmsatzXXV} - 1
  = \Pct{\RdyUmsatzWachstum}
\]

\subsection{Der Fr\"uhindikator: Zulassungsantr\"age und Nordamerika}

Ein Generikahersteller f\"uhrt kein Auftragsbuch; sein Fr\"uhindikator ist die
Zahl der anh\"angigen Zulassungsantr\"age und die Entwicklung des Marktes, in
dem die Margen entstehen. Im Berichtsjahr wurden
\Stueck{\RdyAndaEingereicht} neue Antr\"age eingereicht;
\Stueck{\RdyAndaAnhaengig} sind anh\"angig, kumuliert wurden
\Stueck{\RdyAndaKumuliert} eingereicht.\quelleZF{2026}{55}\quelleMerken{s25zfxcacgxff}
Der nordamerikanische Umsatz des gr\"o{\ss}ten Segments fiel im selben Jahr um
\PctBetrag{\RdyNordamerikaWachstum} auf \MioINR{\RdyNordamerika};
sein Anteil am Segment sank von der H\"alfte auf
\Pct{\RdyNordamerikaAnteil}.\quelleErneut{s25zfxcacgxfe}

\subsection{Rohertragsmarge, ausgewiesen und nach Segment}

Der Rohertrag fiel auf \MioINR{\RdyRohertragXXVI} nach
\MioINR{\RdyRohertragXXV}; die Marge sank von \Pct{\RdyRohertragsmargeVJ} auf
\Pct{\RdyRohertragsmarge}.\quelleErneut{s25zfxcacgxbbc}
\[
  \text{Rohertragsmarge} = \frac{\RdyRohertragXXVI}{\RdyUmsatzXXVI}
  = \Pct{\RdyRohertragsmarge}
\]
Ein Umsatz, der steigt, w\"ahrend der Rohertrag f\"allt, ist der eigentliche
Befund des Jahres. Das Unternehmen benennt die Ursache selbst: Preiserosion
bestehender Produkte einschlie{\ss}lich der Wirkung des
Regalbestandsausgleichs nach der Preissenkung auf
Lenalidomid.\quelleZF{2026}{57}\quelleMerken{s26zfxcacgxfh}
Im gr\"o{\ss}ten Segment sank die Rohertragsmarge von
\Pct{\RdyRohertragGgMargeVJ} auf \Pct{\RdyRohertragGgMarge}.\quelleErneut{s26zfxcacgxfh}
Der bezifferte Abzug aus dem Regalbestandsausgleich betrug
\MioINR{\RdySsaEffekt} und damit \Pct{\RdySsaAnteil} des
Konzernumsatzes.\quelleZF{2026}{158}\quelleMerken{s26zfxcacgxbfi}

\subsection{Operatives Ergebnis}

Das Betriebsergebnis fiel \PctBetrag{\RdyBetriebsergebnisWachstum} auf
\MioINR{\RdyBetriebsergebnis} nach
\MioINR{\RdyBetriebsergebnisVJ}.\quelleZF{2026}{112}\quelleMerken{s26zfxcacgxbbc}
Die operative Marge sank damit von \Pct{\RdyBetriebsMargeVJ} auf
\Pct{\RdyBetriebsMarge}. Enthalten ist eine Wertminderung auf langfristige
Verm\"ogenswerte von \MioINR{\RdyWertminderung}.\quelleErneut{s26zfxcacgxbbc}
Der Forschungsaufwand ging um \PctBetrag{\RdyForschungWachstum} auf
\MioINR{\RdyForschungXXVI} zur\"uck, \Pct{\RdyForschungsquote} des
Umsatzes.\quelleZF{2026}{53}\quelleMerken{s26zfxcacgxfd}

\subsection{Nettogewinn}

Der Jahres\"uberschuss fiel \PctBetrag{\RdyErgebnisWachstum} auf
\MioINR{\RdyJahresueberschussXXVI}.\quelleErneut{s26zfxcacgxbbc}
Je Aktie sind das \INRje{\RdyErgebnisJeAktie} nach
\INRje{\RdyErgebnisJeAktieVJ} im Vorjahr und
\INRje{\RdyErgebnisJeAktieVVJ} im Jahr davor.\quelleErneut{s26zfxcacgxbbc}
Die Eigenkapitalrendite betr\"agt
\[
  \frac{\RdyJahresueberschussXXVI}{\RdyEigenkapital} = \Pct{\RdyEigenkapitalrendite}.
\]
Die Eigenkapitalquote liegt bei \Pct{\RdyEigenkapitalquote}. Den
Finanzschulden von \MioINR{\RdyFinanzschulden} stehen liquide Mittel und
Wertpapiere von zusammen mehr gegen\"uber; das Unternehmen hat eine Nettokasse
von \MioINR{\RdyNettokasse}.\quelleZF{2026}{111}\quelleMerken{s26zfxcacgxbbb}
Die Frage nach der Entschuldung ist damit leer und wird in
Abschnitt~\ref{sec:kapazitaet-rdy} durch die Finanzierungsfrage ersetzt.

\subsection{Dividende und Dividendenrendite}

F\"ur FY2026 wurden \INRje{\RdyDividendeJeAktie} je Aktie ausgesch\"uttet,
unver\"andert gegen\"uber den beiden Vorjahren.\quelleZF{2026}{82}\quelleMerken{s26zfxcacgxic}
Bei einem Kurs von \INRje{\RdyKursRupien} je Aktie ergibt das eine Rendite von
\Pct{\RdyDividendenrendite}; die Aussch\"uttungsquote betr\"agt
\Pct{\RdyAusschuettungsquote}. Eine gleichbleibende Dividende bei fallendem
Gewinn hebt die Quote, ohne dass eine Entscheidung dazu getroffen worden
w\"are.

\subsection{Aktienkursverlauf}
\label{sec:kursbild-rdy}

\begin{figure}[htbp]
\centering
\begin{tikzpicture}
\begin{axis}[kursachserdy]
\addplot+[gray!55, thick, mark=none] table[col sep=comma, x=datum,
  y=kurs] {data/kursverlauf_rdy.csv};
\addplot+[kurspunkte] table[col sep=comma, x=datum,
  y=kurs, meta=art] {data/kursverlauf_rdy.csv};
\end{axis}
\end{tikzpicture}
\caption{Kursverlauf des Hinterlegungsscheins vom September 2016 bis zum
\Datum{\RdyKursTag}, in USD. {\color{kurshoch}$\blacktriangle$} Jahreshoch,
{\color{kurstief}$\blacktriangledown$} Jahrestief,
{\color{kursschluss}$\bullet$} Jahresschluss. Die Reihe ist um den
Aktiensplit vom Oktober 2024 zur\"uckgerechnet.}
\label{fig:kurs-rdy}
\end{figure}

\begin{table}[htbp]
\centering
\caption{Jahresschluss-, H\"ochst- und Tiefstkurse des Hinterlegungsscheins in
USD. Das Jahr 2026 reicht bis zum \Datum{\RdyKursTag}. Sekund\"arquelle: Nasdaq.}
\label{tab:kurs-rdy}
\begin{tabular}{lrrr}
\toprule
Jahr & Schluss & H\"ochst & Tiefst\\
\midrule
\tabellenkoerper{data/kurs_rdy_tabelle}
\bottomrule
\end{tabular}
\end{table}

\FloatBarrier

\subsection*{Kennzahlen\"uberblick}
\addcontentsline{toc}{subsection}{Kennzahlen\"uberblick}

\begin{table}[htbp]
\centering
\caption{Kennzahlen FY2026 gegen FY2025, Betr\"age in Mio.\,INR}
\label{tab:rdy-kennzahlen}
\begin{tabular}{lrrr}
\toprule
Kennzahl & FY2026 & FY2025 & Ver\"anderung\\
\midrule
Umsatz                 & \Betrag{\RdyUmsatzXXVI} & \Betrag{\RdyUmsatzXXV} & \Pct{\RdyUmsatzWachstum}\\
Rohertrag              & \Betrag{\RdyRohertragXXVI} & \Betrag{\RdyRohertragXXV} & \Pct{\RdyRohertragWachstum}\\
Betriebsergebnis       & \Betrag{\RdyBetriebsergebnis} & \Betrag{\RdyBetriebsergebnisVJ} & \Pct{\RdyBetriebsergebnisWachstum}\\
Jahres\"uberschuss     & \Betrag{\RdyJahresueberschussXXVI} & \Betrag{\RdyJahresueberschussXXV} & \Pct{\RdyErgebnisWachstum}\\
Ergebnis je Aktie (INR) & \Kurswert{\RdyErgebnisJeAktie} & \Kurswert{\RdyErgebnisJeAktieVJ} & \Pct{\RdyEpsWachstum}\\
Forschungsaufwand      & \Betrag{\RdyForschungXXVI} & \Betrag{\RdyForschungXXV} & \Pct{\RdyForschungWachstum}\\
Freier Mittelzufluss   & \Betrag{\RdyFcfXXVI} & \Betrag{\RdyFcfXXV} & \Pct{\RdyFcfWachstum}\\
Eigenkapital           & \Betrag{\RdyEigenkapital} & \Betrag{\RdyEigenkapitalVJ} & \Pct{\RdyEigenkapitalWachstum}\\
\bottomrule
\end{tabular}
\end{table}

\subsection{Die Mehrjahresreihe und die Lage im Zyklus}
\label{sec:reihe-rdy}

\begin{table}[htbp]
\centering
\caption{Neun Jahre, Betr\"age in Mio.\,INR; freier Mittelzufluss je Aktie in
INR, durchg\"angig auf die Aktienzahl nach dem Split bezogen}
\label{tab:rdy-reihe}
\begin{tabular}{lrrrrr}
\toprule
Jahr & Umsatz & Jahres\-\"uber\-schuss & operativer
     & freier & je Aktie\\
     &        &                        & Mittelzufluss & Mittelzufluss & \\
\midrule
\tabellenkoerper{data/flow_rdy_reihevoll}
\bottomrule
\end{tabular}
\end{table}

Der freie Mittelzufluss ist hier operativer Mittelzufluss abz\"uglich der
Investitionen in Sachanlagen und immaterielle Verm\"ogenswerte zuz\"uglich der
Erl\"ose aus deren Abgang. Die Reihe schwankt stark: zwischen
\INRje{\RdyFcfJeAktieXVIII} im Jahr FY2018 und dem Vierfachen davon in FY2023.

Zwei Ablesungen. Erstens liegt FY2026 mit \INRje{\RdyFcfVoll} nahe dem Median
der neun Jahre (\INRje{\RdyFcfMedian}). Anders als bei AstraZeneca ist das
Basisjahr hier kein Spitzenwert, sondern Mittelfeld. Zweitens h\"angt die
Wachstumsrate vollst\"andig vom gew\"ahlten Startjahr ab: von FY2018
gerechnet sind es \Pct{\RdyCagr} pro Jahr, von FY2019 gerechnet
\Pct{\RdyCagrSieben}. Dieselbe Reihe, zwei Vorzeichen.
Abschnitt~\ref{sec:verdikt-rdy} macht daraus keine Wahl, sondern einen Befund.
